\documentclass[11pt]{article}
\usepackage{amsmath,amssymb,amsfonts}
\usepackage{geometry}
\usepackage{enumitem}
\usepackage{color}
\usepackage{hyperref}
\usepackage{booktabs}

\geometry{letterpaper, margin=1in}

\definecolor{addressed}{rgb}{0,0.5,0}
\definecolor{partial}{rgb}{0.8,0.4,0}
\definecolor{pending}{rgb}{0.8,0,0}

\title{\textbf{Response to Reviewers}\\[0.5em]
\large Manuscript: Flux-Preserving Adaptive Finite State Projection\\
for Stiff Stochastic Systems\\[0.5em]
\normalsize Submitted to: SIAM Multiscale Modeling and Simulation}

\author{Aditya Dendukuri, Shivkumar Chandrasekaran, and Linda Petzold\\
Department of Computer Science\\
University of California, Santa Barbara}

\date{\today}

\begin{document}

\maketitle

\noindent Dear Editor and Reviewers,

We thank the reviewers for their careful and constructive evaluation of our manuscript.
In light of the reports, we have undertaken a major revision that changes both the
scope and the technical core of the paper.

The original submission focused on a \emph{quantile-based} pruning strategy: at each
time step, a fixed fraction of probability mass was removed and the distribution was
renormalized. As Reviewer~1 rightly pointed out, such probability-only pruning can
fail badly in the presence of bottleneck reactions, and the error analysis in that
version was not rigorous enough.

In the revised manuscript we now present a \emph{flux-preserving adaptive FSP}
framework whose central ideas are:

\begin{enumerate}[topsep=0pt,itemsep=2pt]
  \item Use \emph{probability flux}---rather than probability alone---to identify
        states that are dynamically important but may be rarely occupied
        (Definitions~3.1–3.2).
  \item Combine quantile-based pruning with a \emph{flux safeguard} that prevents
        the removal of low-probability, high-flux bottleneck states (Section~3.2
        and Algorithm~3.1).
  \item Drive the time-step selection by the \emph{total system flux}, leading to a
        flux-adaptive time-step controller that automatically shortens the time
        step during stiff or transient phases and enlarges it in slow regimes
        (Corollary~4.8).
  \item Derive a new error analysis based on a block-decomposed generator and
        boundary fluxes, yielding rigorous local and global $\ell_1$ error bounds
        via contraction of the stochastic semigroup (Section~4, especially
        Proposition~4.1, Theorem~4.4, and Corollaries~4.6–4.8).
  \item Demonstrate the algorithm on demanding benchmark problems, including a
        bottleneck system, a stochastic toggle switch, the oscillatory Oregonator,
        and the stiff Robertson (ROBER) autocatalytic system, with detailed
        reporting of algorithmic parameters, state-space sizes, and performance
        (Sections~3.3, 4.3, and~5).
\end{enumerate}

Below we respond point-by-point to the reviewers' comments.

%======================================================================
\section*{Response to Reviewer 1}
%======================================================================

\subsection*{Major Comments}

\noindent\textbf{Comment 1.1:} \textit{Lines 178–180: Can the authors elaborate why using a fixed probability
cutoff ``can be unreliable when the overall probability distribution changes
drastically over time''?}

\textbf{Response:} We have considerably expanded both the explanation and the
numerical evidence for this point.

Conceptually, the difficulty is that low probability does not imply negligible
dynamical importance: states with small $p(x,t)$ but large exit rates $w(x)$ can
carry significant probability flux and act as essential conduits between regions of
state space. This is now formalized in Definitions~3.1–3.2, where we define the
outgoing probability flux
\[
  \Phi(x,t) = p(x,t)\,w(x)
\]
and the total system flux
\[
  \Phi_{\mathrm{total}}(J,t)
  = \sum_{x\in J} \Phi(x,t),
\]
and we explicitly point out that such states can be connectivity-critical despite
being rarely occupied (lines 148–156).

To make this concrete, the revised Introduction and Section~3.3 now use the
reviewer’s bottleneck idea:
\[
  A \xrightarrow{k_1} B, \qquad B \xrightarrow{k_2} B + C,
\]
with $k_1 \ll k_2$ and initial condition $(N_A,N_B,N_C) = (1,0,0)$. In this system
the intermediate state with $(N_A,N_B)=(0,1)$ maintains very small probability but
large exit flux, and it is the \emph{only} pathway to increasing $C$.

Section~3.3 shows (Figure~3 and Table~1) that:

\begin{itemize}
  \item Under the original, probability-only pruning (quantile step without flux
        safeguard), states on this bottleneck path are removed as soon as their
        probability falls below the quantile threshold. This makes the truncated
        generator reducible and effectively traps probability near $(1,0,0)$.
  \item Under the new flux-preserving algorithm, the same states are protected
        because their $\Phi(x,t)$ constitutes a non-negligible fraction of the
        total system flux, even though $p(x,t)$ is tiny. As a result, the
        algorithm correctly captures the transition to large $C$ and the advancing
        probability wave front.
\end{itemize}

Thus the revised manuscript both explains and demonstrates why fixed probability
cutoffs can be unreliable and how the flux safeguard resolves the failure mode.

\vspace{0.5em}

\noindent\textbf{Comment 1.2:} \textit{Is it necessary to always renormalize the generator matrix and the
probability distribution? One powerful property of FSP comes from controlling
the loss of probability mass via $1-\mathbf{1}^\top\mathbf{p}(t)$. Also, if the
truncated generator is simply the principal submatrix of the full generator and
is never renormalized, then $1-\mathbf{1}^\top\mathbf{p}(t)$ is a global error
estimate. Can the authors elaborate why normalization is favorable or necessary
and why their more elaborate error control is better?}

\textbf{Response:} We thank the reviewer for highlighting this important point. The
revised manuscript clarifies the distinction between \emph{generator} normalization
and \emph{probability} normalization, and explains how our flux-based error control
relates to the classical FSP bound.

\begin{itemize}
  \item \emph{Generator:} In the revised method we \emph{do not} renormalize the
        generator. Instead, Section~4.1 constructs a \emph{compressed generator}
        $A_{JJ}$ on the active set $J$ by folding transitions into and out of
        the complement $J'$ into a reflecting boundary term (equation~(4.4)). As
        shown in Proposition~4.1, this compressed generator satisfies
        $\mathbf{1}^\top A_{JJ} = 0$ and hence generates a contraction
        semigroup in $\ell_1$. This is conceptually similar to taking a principal
        submatrix but retains the effect of the truncated region through the
        boundary term.

  \item \emph{Probability:} The probability vector on $J$ is renormalized after
        each pruning event so that $\|\mathbf{p}_J(t_n)\|_1 = 1$ holds at the
        discrete times $t_n$. This keeps the numerical solution interpretable as
        a probability distribution on the current active set and improves
        numerical robustness when the remaining mass is small. In the error
        analysis, the renormalization contribution is explicitly accounted for in
        the local model error and is bounded by the pruned mass.

  \item \emph{Relation to $1-\mathbf{1}^\top\mathbf{p}_J(t)$:} For a fixed FSP
        with generator equal to a principal submatrix of the full generator,
        $1-\mathbf{1}^\top\mathbf{p}_J(t)$ indeed bounds the probability that has
        left $J$. Our block-partitioned analysis in Section~4 recovers this view in
        more detail. In Theorem~4.4 we express the local model error over a time
        step $[t_n,t_{n+1}]$ in terms of the \emph{boundary fluxes}
        $\Phi_{\mathrm{out}}$ and $\Phi_{\mathrm{in}}$ between $J$ and $J'$. Then
        Corollary~4.6 shows that the global error after $N$ steps satisfies
        \[
          e_N \le \sum_{n=0}^{N-1} (\varepsilon_n + \tau_n),
        \]
        where $\varepsilon_n$ is an integral of these boundary fluxes and $\tau_n$
        is the numerical time-stepping error.
\end{itemize}

Compared to a single scalar quantity $1-\mathbf{1}^\top\mathbf{p}_J(t)$, the
flux-based decomposition distinguishes (i) model error due to truncation and
pruning (through boundary fluxes and pruned mass) from (ii) numerical error due to
the matrix exponential. This structure directly motivates our adaptive strategy,
which uses the boundary fluxes to decide when to expand the state space and how to
adjust the time step. We have added a short remark in Section~4.2 to make this
connection explicit.

\vspace{0.5em}

\noindent\textbf{Comment 1.3:} \textit{The proof for Theorem 3.10 is incomplete. Specifically, I do not
understand why the errors could simply be summed over the $N$ time steps.}

\textbf{Response:} We agree that the previous argument was incomplete. In the revised
manuscript, the old Theorem~3.10 has been removed and replaced by a new error
analysis in Section~4, culminating in Corollary~4.6 (global error bound).

The revised analysis proceeds as follows:

\begin{itemize}
  \item Proposition~4.1 proves that the compressed generator $A_{JJ}$ generates a
        contraction semigroup in $\ell_1$, i.e.,
        $\|\mathrm{e}^{A_{JJ} t}\|_{1\to 1} = 1$ for $t\ge 0$.
  \item Theorem~4.4 bounds the \emph{local model error} over one time step in terms
        of the integrated boundary fluxes across the truncation.
  \item Proposition~4.5 bounds the \emph{local time-stepping error} for the
        Krylov-based matrix exponential approximation, controlled by the ODE
        tolerance $\varepsilon_{\mathrm{ODE}}$.
  \item Corollary~4.6 then combines these pieces to show that
        \[
          e_{n+1} \le e_n + \varepsilon_n + \tau_n,
        \]
        and by induction
        \[
          e_N \le \sum_{n=0}^{N-1} (\varepsilon_n + \tau_n),
        \]
        where $e_n$ is the global $\ell_1$ error at time $t_n$.
\end{itemize}

Because the semigroup is contractive, local errors do \emph{not} get amplified as
they are propagated forward in time, which justifies the additive accumulation in
Corollary~4.6. We now make this point explicit in the proof.

\vspace{0.5em}

\noindent\textbf{Comment 1.4:} \textit{Certain fringe models with “bottleneck reactions” might be nontrivial
with any FSP with time-stepping or sliding windows, including the method in this
paper. [A specific three-species bottleneck system is described and the reviewer
asks the authors to comment.]}

\textbf{Response:} We are very grateful for this comment; it directly motivated one of
the central case studies in the revised manuscript.

Section~3.3 now presents a detailed bottleneck example closely following the
reviewer’s suggestion. We consider a three-species system with a slow activation step
and a fast production step leading to an advancing front in $C$. The parameters are
chosen so that the intermediate states on the activation path have extremely low
probability but carry the only significant flux into regions with large $C$.

In this setting we compare:

\begin{itemize}
  \item a probability-only variant (quantile pruning without flux safeguard), and
  \item the full flux-preserving algorithm (Algorithm~3.1 combined with the
        flux-adaptive time stepping of Section~3.2).
\end{itemize}

Figure~3 and Table~1 show that the probability-only variant quickly removes
bottleneck states from the truncated set, after which the numerical solution remains
localized near the initial state and fails to reproduce the growth of $C$. In
contrast, the flux-preserving variant retains bottleneck states whenever their exit
flux remains a non-negligible fraction of the total system flux, and it reproduces
the expected increase in $\langle C\rangle$ over several orders of magnitude in time.
The example thus illustrates precisely the kind of pathology the reviewer described
and demonstrates that the flux safeguard is effective in protecting such bottleneck
channels.

\vspace{0.5em}

\noindent\textbf{Comment 1.5:} \textit{Can the authors compare their method to previous FSP methods that
also incorporate adaptive state spaces, either in discussion or numerical tests?
Some are already cited in the paper. Another paper that comes to mind is Dinh \&
Sidje (Physical Biology, 2017).}

\textbf{Response:} We have expanded the discussion of related work in the
Introduction and in Section~2.3 to position the revised method more clearly among
existing approaches.

In particular, we now discuss:

\begin{itemize}
  \item \emph{Aggregation and lumping methods} such as those of Bobbio \& Trivedi and
        Zhang, Watson, \& Cao, which group micro-states into macro-states using
        stationary or Monte Carlo information. These methods achieve substantial
        reduction but rely on either a priori timescale separation or sufficiently
        rich sampling.
  \item \emph{Sliding-window and hybrid FSP–SSA schemes} such as those by Wolf
        et al.\ and MacNamara et al., which track moving regions of support or
        combine deterministic and stochastic updates.
  \item \emph{Krylov-FSP-SSA and tensor-train methods} of Dinh \& Sidje and
        collaborators, which guide state-space expansion using SSA trajectories or
        compressed representations (as cited in references [6–7]).
\end{itemize}

Our flux-preserving adaptive FSP differs in two main ways:

\begin{enumerate}
  \item We use \emph{absolute boundary flux} as the main control quantity, which
        leads directly to the local and global error bounds in Section~4, rather
        than relying on relative flux thresholds or stationary information.
  \item We decouple discovery of new states (via reachability expansion of the
        reaction network) from the pruning decision, using quantiles only as a
        pre-selection mechanic and flux as the final safeguard. This explicitly
        addresses the bottleneck scenario highlighted by the reviewer.
\end{enumerate}

Implementing and tuning several competing adaptive FSP codes for a full numerical
benchmark on all our examples would require a substantial additional effort and is
beyond the scope of the current revision. Instead we focus on clearly articulating
the conceptual differences, and we hope the detailed algorithmic description and
performance data in the revised manuscript will facilitate future comparisons by
other groups.

\vspace{0.5em}

\noindent\textbf{Comment 1.6:} \textit{Although well studied, Michaelis–Menten, Lotka–Volterra, and toggle
switch models are relatively simple to solve. It would be beneficial to showcase
your algorithm with at least one example where the state space is known to be
large. Can the authors implement such a model?}

\textbf{Response:} We agree that more demanding examples are needed to demonstrate the
advantages of the method. The revised manuscript contains several such systems.

\begin{itemize}
  \item Section~5.1 revisits the stochastic toggle switch, but now primarily as a
        qualitative illustration of bistability and basin switching under the
        flux-preserving truncation (Figure~6).
  \item Section~5.2 introduces the \textbf{Oregonator} model of chemical
        oscillations. The natural bounding box for the molecular copy numbers has
        on the order of $3\times 10^7$ states, whereas the adaptive FSP keeps
        between a few hundred and a few thousand states active at any time,
        achieving a compression ratio of roughly $10^4:1$ while faithfully
        reproducing the oscillatory dynamics (Figure~7).
  \item Section~5.3 adds the \textbf{Robertson (ROBER) autocatalytic system}, which
        is extremely stiff, with reaction rates spanning nine orders of magnitude.
        A conservative rectangular bounding box guided by the mean trajectories
        would contain on the order of $10^8$–$10^9$ states, whereas the adaptive
        state space in our simulations remains in the range $50$--$2000$ states.
        Thus the method achieves a compression of several orders of magnitude
        while still capturing the three-phase dynamics (slow initiation, explosive
        autocatalysis, and long equilibration) and controlling local boundary flux
        and time-stepping error (Figure~8 and Table~5).

\end{itemize}

These examples go beyond the original low-dimensional tests and show that the
flux-preserving adaptive FSP can handle genuinely large and stiff systems.

\subsection*{Minor Comments}

\begin{enumerate}
  \item \textbf{Line 148:} The extra period has been removed in the revised text.
  \item \textbf{Line 163:} The colon has been replaced by a period.
  \item \textbf{Algorithm 2.1:} The termination condition has been corrected so that
        it reflects the intended tolerance on the missing probability mass.
  \item \textbf{Lines 178–179:} We now explicitly say that the ``total probability''
        refers to the cumulative probability used in the Optimal FSP truncation
        criterion.
\end{enumerate}

%======================================================================
\section*{Response to Reviewer 2}
%======================================================================

\subsection*{Main Comments}

\noindent\textbf{Comment 2.1:} \textit{The first experiment, depicted in Figure 2, requires a much deeper
discussion. The error seems to increase with time, contradicting the analysis.
Also, 1000 SSA trajectories would not be enough to eliminate statistical error.}

\textbf{Response:} We agree that the original SSA-based experiment did not clearly
separate truncation error from Monte Carlo sampling error, and that the discussion
of error growth was incomplete.

In the revised manuscript we have \emph{replaced} this experiment by a deterministic
error study that directly tests the theoretical results:

\begin{itemize}
  \item Section~4.3 (\emph{Experimental Error Analysis}) compares the
        flux-preserving adaptive FSP to a high-accuracy fixed-FSP reference on a
        bounded state space $\{0,\dots,2\}\times\{0,\dots,2\}\times\{0,\dots,10000\}$.
  \item Figure~5 shows that the relative error in $\langle C\rangle$ remains below
        $0.1\%$ and decreases to about $0.015\%$ as the distribution spreads
        (panel~5a). The instantaneous error rate is approximately constant at
        $4.9\times 10^{-7}\,\mathrm{s}^{-1}$ (panel~5b), in agreement with the
        additive error structure of Corollary~4.6. Panel~5d confirms linear
        scaling of the absolute error with $\langle C\rangle$, corresponding to
        about $0.01\%$ relative error over three orders of magnitude in
        population.
  \item Table~2 reports final-time means and errors for all species at $t=10^5$~s,
        with the adaptive and reference solutions in close agreement and the
        adaptive state space about $4.5$ times smaller than the fixed one.
\end{itemize}

We believe this new experiment more directly and transparently validates the
revised error analysis.

\vspace{0.5em}

\noindent\textbf{Comment 2.2:} \textit{For most (if not all) experiments, the pruning factor $\alpha$, the
effective pruned fraction $m$, and the time step are not reported. These should
be specified in the paper rather than needing to be inferred from the code.}

\textbf{Response:} We have revised the numerical sections to state the key algorithmic
parameters explicitly:

\begin{itemize}
  \item Section~3.3, which introduces the flux-preserving pruning algorithm, now
        specifies the quantile tolerance $\alpha$, flux tolerance $\varepsilon_{\rm
        flux}$, and the bounds $\Delta t_{\min}$ and $\Delta t_{\max}$ used in the
        bottleneck experiments, as well as the functional form of the flux-based
        time-step rule.
  \item Section~4.3 provides the time grid and state-space bounds for the fixed-FSP
        reference solution, and the adaptive time-step settings for the comparison
        run.
  \item Sections~5.1–5.3 state the choices of $\alpha$, flux tolerance, and
        time-step range used in the toggle, Oregonator, and ROBER examples. The
        figures also show the resulting adaptive time steps and state-space sizes
        as functions of time.
\end{itemize}

In the revised method, $\alpha$ primarily determines how much probability mass is
\emph{eligible} for pruning per step, while the flux safeguard further constrains
which of those states can actually be removed. Rather than reporting an effective
$m$ at every time step, we focus on how $\alpha$, the flux threshold, and the
boundary fluxes enter the local and global error bounds in Section~4.

\vspace{0.5em}

\noindent\textbf{Comment 2.3:} \textit{I would like to see more explicitly the relation between the pruning
factor $\alpha$ and the effective pruned fraction $m$. This relationship will
affect both efficiency and accuracy.}

\textbf{Response:} We appreciate this request and have clarified the role of $\alpha$
in the revised algorithm.

In the original quantile-only scheme, $\alpha$ directly controlled the fraction of
probability mass removed at each pruning step, so $m \approx \alpha$. The new
method introduces an additional flux-based safeguard:

\begin{itemize}
  \item The quantile parameter $\alpha$ selects a \emph{candidate set} of low-probability
        states whose cumulative probability does not exceed $\alpha$.
  \item The flux safeguard then scans this candidate set and retains any state
        whose outgoing flux exceeds a prescribed fraction of the total system flux.
        These protected states are not pruned, regardless of their probability.
\end{itemize}

As a consequence, the \emph{effective} pruned mass $m$ is at most $\alpha$, and can
be considerably smaller when bottleneck states are present. This behavior is
illustrated in the bottleneck experiment: even with a relatively large $\alpha$, the
flux safeguard prevents removal of the intermediate bottleneck states and thereby
limits $m$ while maintaining connectivity. Equation~(4.22) and the discussion in
Section~4.2 make explicit how the local model error depends on the pruned mass, the
boundary fluxes, and the time step, thereby connecting $\alpha$ and $m$ through
their impact on the flux-based error control.

A full parametric study of $(\alpha,m)$ over many systems would significantly
lengthen the paper; we therefore leave a more systematic investigation of this
relationship to future work.

\vspace{0.5em}

\noindent\textbf{Comment 2.4:} \textit{The authors report runtimes for their new algorithm, but these are
only average values on a single hardware platform. At least the hardware used
should be mentioned, and runtimes should be compared to those of existing
algorithms; probably higher-dimensional examples need to be studied as well.}

\textbf{Response:} We agree that performance information and reproducibility are
important. In the revised manuscript we have:

\begin{itemize}
    \item Specified the implementation details and hardware in Section~4.1: all
        experiments were run in Julia using Expokit.jl for Krylov-based matrix
        exponential computations, Catalyst.jl for reaction-network modeling, and
        sparse matrix structures from the \texttt{SparseArrays} standard library.
        Simulations were performed on a MacBook Pro with an M1 Pro chip and
        32\,GB of RAM.
  \item Added Section~6 and Section~6.1, which analyze the performance of the
        \emph{matrix reconstruction} strategy that reuses generator entries when
        consecutive active state spaces overlap. Figure~9 and Table~6 show that,
        for the Oregonator model, reconstruction reduces the median time for
        generator construction from about $19.2$~ms to $1.25$~ms, a speedup of
        roughly $15\times$ at the kernel level. Figure~10 and Table~7 then report
        total simulation times, where matrix reconstruction yields speedups of
        about $3.9\times$ (Oregonator), $3.5\times$ (ROBER), and $2.9\times$
        (toggle switch) relative to a standard implementation that rebuilds the
        matrix from scratch at each time step.
\end{itemize}

A direct head-to-head comparison with all existing adaptive FSP implementations
(e.g., Optimal FSP, Krylov-FSP-SSA, aggregation-based methods) on the same
hardware would require reimplementing or interfacing with multiple external codes,
and carefully tuning each method for fairness. We believe this is an interesting
direction for future work but beyond the scope of the present revision. Our goal
here is to document the internal performance characteristics of the proposed
algorithm and to provide enough detail for readers to compare against their own
implementations.

\vspace{0.5em}

\noindent\textbf{Comment 2.5:} \textit{Include a test case where excessive pruning would be an issue but
your new method avoids that problem.}

\textbf{Response:} The bottleneck example in Section~3.3 was designed precisely to
address this request and the related concerns raised by Reviewer~1, who suggested
this type of system in their report. This example became the starting point for
shifting the scope of the paper toward stiff, multiscale systems where rare but
dynamically essential transitions must be preserved.

As described in our response to Comment~1.4, the bottleneck model contains
low-probability states that form the only pathway to large $C$. Any scheme that
prunes solely on probability is forced either to choose a very small $\alpha$
(sacrificing efficiency) or eventually to remove these states and thus
qualitatively change the dynamics. The flux-preserving adaptive FSP method, by
contrast, monitors the outgoing flux from candidate states and protects those whose
flux is non-negligible. The plots and tables in Section~3.3 show that the
flux-based variant retains the bottleneck channel and recovers the correct
long-time behavior, whereas the probability-only variant does not.


\subsection*{Minor Comments}

\noindent\textbf{Comment:} \textit{The last sentence of the abstract is not a sentence.}

\textbf{Response:} We have rewritten the abstract to consist of complete sentences and
to reflect the revised focus on flux-based pruning and error control.

%======================================================================
\section*{Summary}
%======================================================================

In summary, the revised manuscript differs substantially from the original
submission. The main advances are:

\begin{enumerate}
  \item A \emph{flux-preserving} adaptive FSP algorithm that augments
        quantile-based pruning with a flux safeguard and a flux-guided time-step
        controller, directly addressing the bottleneck failure modes pointed out
        by Reviewer~1.
  \item A new \emph{error analysis} that expresses local model error in terms of
        boundary fluxes and exploits semigroup contraction to obtain additive
        global error bounds (Theorem~4.4, Proposition~4.5, and
        Corollaries~4.6–4.8).
  \item A set of \emph{challenging numerical experiments}, including a detailed
        bottleneck case study, a stochastic toggle switch, the oscillatory
        Oregonator, and the stiff Robertson system, each with explicit reporting of
        parameters, state-space sizes, and observed behavior.
  \item A \emph{performance study} of matrix reconstruction that demonstrates
        substantial speedups for large, slowly evolving state spaces (Section~6).
\end{enumerate}

We believe these changes fully address the reviewers' concerns and substantially
strengthen the manuscript. We are very grateful to both reviewers and to the
Associate Editor for their guidance, which has been instrumental in improving this
work.

\vspace{1em}
\noindent Sincerely,\\[0.5em]
Aditya Dendukuri, Shivkumar Chandrasekaran, and Linda Petzold


\end{document}
